\chapter{Considerações Finais}\label{cap:conclusao}

A \gls{ca} consolida-se como um paradigma essencial para suprir a crescente demanda por desempenho computacional e eficiência energética, especialmente diante das limitações físicas na evolução do \textit{hardware}. Técnicas de aproximação foram propostas para suprir essa demanda, como é o caso de: perforação de laços, a memoização aproximada e o relaxamento em operações de ponto flutuante. Este trabalho estas técnicas foram introduzidas como uma extensão para a \gls{api} \texttt{OpenMP} na infraestrutura do LLVM, o objetivo foi facilitar a integração e o controle desses algoritmos aproximados em aplicações paralelas.

A validação ocorreu por meio de experimentos computacionais envolvendo sete aplicações distintas, nas quais foram analisados o tempo total de execução e as métricas de qualidade dos resultados gerados. De maneira geral, a cláusula \texttt{perfo\_large}, que aplica a perforação de laços modificando diretamente a variável de indução, obteve os ganhos de desempenho mais expressivos, reduzindo significativamente o tempo de execução na maioria dos cenários, ainda que ao custo de uma degradação na acurácia das saídas.

Entre as demais abordagens, as estratégias \texttt{perfo\_init} e \texttt{perfo\_fini} atuaram como alternativas intermediárias, entregando ganhos de desempenho moderados, porém inferiores aos alcançados pela \texttt{perfo\_large}. A cláusula \texttt{memo} demonstrou um bom potencial de aceleração, mas restrito a cenários específicos com alta previsibilidade de entrada. Por fim, a diretiva \texttt{fastmath} não superou o desempenho da execução exata padrão.

Como trabalhos futuros, sugere-se estender o suporte dessas cláusulas de aproximação para o nível de GPU, o que permitiria explorar o paralelismo massivo desses aceleradores em conjunto com os benefícios de desempenho da computação aproximada.